%! TEX program = lualatex
\documentclass[10pt, a4paper]{beamer} %handout fuer eine gedruckte version der slides
\usetheme{metropolis}

\usepackage{url}
\usepackage{polyglossia}
\setmainlanguage{english}

\usepackage{listings,xcolor}
\usepackage{graphicx}
\usepackage{nameref}
\usepackage{mathtools}

\usepackage{booktabs}

\makeatletter
\newcommand*{\currentname}{\@currentlabelname}
\makeatother

\usepackage{fontspec}
% \setmainfont{UbuntuMono}
\newfontfamily\Bera{Bitstream Vera Sans Mono}[Scale=0.85]

\definecolor{mDarkTeal}{HTML}{23373b}
\definecolor{mLightBrown}{HTML}{EB811B}
\definecolor{mDarkBlue}{HTML}{0F3470}
\definecolor{mDarkGrey}{HTML}{999999}
\definecolor{mLighterGrey}{HTML}{CCCCCC}
\definecolor{mLightGrey}{rgb}{0.95,0.95,0.95}
\definecolor{mLightBlue}{HTML}{212751}

\newcommand{\lb}[1]{{\color{mLightBrown}#1}}

\newcommand{\bblock}[3]{
        \begin{block}{#1}
        #2
    \end{block}
}

\setbeamercolor{block title}{fg=mDarkTeal}

\lstset{language=Python,
    backgroundcolor=\color{mLightGrey},
    frame=single,
    rulecolor=\color{mLightGrey}, % make frame "invisible"
    basicstyle=\Bera\footnotesize,
    keywordstyle=\bfseries\color{mDarkBlue},
    commentstyle=\color{mDarkGrey}\itshape,
    captionpos=t,
    emphstyle=\ttb\color{mDarkBlue},    % Custom highlighting style
    stringstyle=\color{mLightBrown},
    tabsize=2,
    numberbychapter=false,
    showstringspaces=false,
    breaklines=true,
    morekeywords={__init__, and, assert, break, class, continue, def, del, elif, else, except, exec, finally, for, from, global, if, import, in, is, lambda, not, or, pass, print, raise, return, try, while, yield}
}
%
\makeatletter
\setbeamertemplate{footline}{%
\leavevmode
\vbox{\begin{beamercolorbox}[dp=1.25ex,ht=2.75ex]{fg=black}%
  \hspace*{1em}\insertsectionhead%
  \ifx\insertsubsectionhead\@empty\relax\else\hspace*{\fill}\insertsubsectionhead\hspace*{2ex}\fi
  \end{beamercolorbox}%
  }%
}
\makeatother

\setbeamertemplate{section in toc}{%
  {\color{mDarkTeal}\rule[0.3ex]{3pt}{3pt}}~\inserttocsection\par}
\setbeamertemplate{subsection in toc}{%
  \hspace{1.2em}{\color{mLightBrown}\rule[0.3ex]{3pt}{3pt}}~\inserttocsubsection\par}


% \usecolortheme{spruce}
\AtBeginSection[]
{
  \begin{frame}
    \frametitle{\currentname}
    \tableofcontents[currentsection]
  \end{frame}
}

\AtBeginSubsection[]
{
  {
  % \setbeamercolor{background canvas}{bg=mLighterGrey}
  \begin{frame}
    \frametitle{\currentname}
    \tableofcontents[currentsubsection]
  \end{frame}
  }
}

\title % (optional, only for long titles)
{Exercises}
\author % (optional, for multiple authors)
{Philipp Gloor\inst{1}}
\institute
{
  \inst{1}%
  University of Zurich
}
\date{}
\subject{Python}
\titlegraphic{\includegraphics[width=0.3\textwidth]{pics/uzh_logo_e_pos}}
\begin{document}
\begin{frame}
\titlepage
\end{frame}

{
\setbeamercolor{background canvas}{bg=mLightBlue}
\setbeamercolor{frametitle}{fg=white,bg=mLightBlue}
\setbeamercolor{normal text}{fg=white}
\usebeamercolor[fg]{normal text}
\bfseries
\begin{frame}[c, fragile, allowframebreaks]\frametitle{Exercise 1}
    
\begin{itemize}
    \item Write a function \texttt{compare(x,y)} that
    \begin{itemize}
        \item prints 1 if $x > y$
        \item prints 0 if $x == y$
        \item prints -1 if $x < y$
    \end{itemize}
    \item Use \texttt{input()} to receive user input
\end{itemize}

\framebreak

\begin{itemize}
    \item Attention: \texttt{input()} stores the input as a string (not as a number)
    \item If the input is supposed to be a number (int, float) you need to convert it
\end{itemize}

{
\mdseries
\setbeamercolor{normal text}{fg=black}
\usebeamercolor[fg]{normal text}
\begin{lstlisting}
first_number = input('Please enter a first number ')
first_number = int(first_number)
second_number = input('Please enter a first number ')
second_number = int(second_number)
result = first_number + second_number
print(str(result))
\end{lstlisting}
}

\end{frame}

\begin{frame}[c, fragile]\frametitle{Exercise 2}
    
\begin{itemize}
    \item Write a function called \texttt{distance(x1, y1, z1, x2, y2, z2)} which computes the distance between point 1 $(x1, y1, z1)$ and point 2 $(x2, y2, z2)$
    \item Note:
        \begin{itemize}
            \item \( \text{distance} = \sqrt{\left(x_2 - x_1\right)^2 + \left(y_2 - y_1\right)^2 + \left(z_2 - z_1\right)^2} \)
            \item $x^2$ is represented by \texttt{x**2} in Python
            \item The root of x is computed with math.sqrt(x)
            \item Use the import math statement at the beginning of the file
        \end{itemize}
\end{itemize}
\end{frame}

\begin{frame}[c, fragile]\frametitle{Exercise 3}
\begin{itemize}
    \item Write a function \texttt{volume\_from\_radius(radius)}, which calculates the
volume of a sphere
    \item Note:
    \begin{itemize}
        \item \( \text{volume} = \frac{4\pi}{3}\cdot r^3\)
        \item Pi is math.pi
    \end{itemize}
\end{itemize}
    
\end{frame}


\begin{frame}[c, fragile]\frametitle{Exercise 4}
\begin{itemize}
    \item Write a function \texttt{volume\_from\_points(x1, y1, z1, x2, y2, z2)}
    \item This function calculates the volume of a sphere whose radius is the distance between the points $(x1, y1, z1)$ and $(x2, y2, z2)$
    \item Tip: Use the implemented methods from the previous exercises (you might want to copy and paste them into a new file)
\end{itemize}
    
\end{frame}

\begin{frame}[t, fragile]\frametitle{Exercise 5}
\begin{itemize}
    \item Write a function \texttt{is\_between(x, y, z)} which returns True if x <= y <=
z and False otherwise
\end{itemize}
    
\end{frame}

\begin{frame}[c, fragile]\frametitle{Exercise 6}
    
\begin{itemize}
    \item Write a function \texttt{calc\_sum(numbers)} which expects a list of numbers as input and returns their sum
    \item The method should be called as follows \texttt{calc\_sum([4,6,7])}
\end{itemize}
\end{frame}

\begin{frame}[c, fragile]\frametitle{Exercise 7}
    
\begin{itemize}
    \item Write a function \texttt{print\_reverse(text)} which expects a string as an argument and prints every character of the string in reverse order
    \item Use a while loop to do this
\end{itemize}
\end{frame}

\begin{frame}[c, fragile]\frametitle{Exercise 8}
\begin{itemize}
\item Write a function \texttt{count\_words(words, min\_word\_length)} that counts the number of words in a list that are at least as long as the specified word length
\item Use a for loop to do this
\item Example:
\end{itemize}

{
\mdseries
\setbeamercolor{normal text}{fg=black}
\usebeamercolor[fg]{normal text}
\begin{lstlisting}
words = ['Emanuel', 'John', 'Ale']
count_words(words, 4)
# 2
\end{lstlisting}
}

\end{frame}

\begin{frame}[c, fragile, allowframebreaks]\frametitle{Exercise 9}
    
\begin{itemize}
    \item Write a function \texttt{calculate\_mark(points, max\_points)} which returns a grade in the Swiss grading scale
    \[ \text{mark} = \frac{\text{points}\cdot 5}{\text{max points}} + 1 \]
    \item The function rounds the grade to the nearest 0.5
    \begin{itemize}
        \item 5.66666 -> 5.5
        \item 5.75 -> 6
    \end{itemize}
    \item The function should accept strings as arguments
    \item Arguments should be therefore converted to floats
\end{itemize}

\framebreak

\begin{itemize}
    \item Write code that asks for points and max\_points as long as the user does not enter "exit"
    \item The grade should be printed after each run
\end{itemize}

{
\mdseries
\setbeamercolor{normal text}{fg=black}
\usebeamercolor[fg]{normal text}
\begin{lstlisting}
# input max_points
while True:
    # input points (use input)
    if points == 'exit':
        break
    # call calculate_mark function
    # print result
\end{lstlisting}
}
\framebreak

\begin{itemize}
    \item Change your code that it additionally asks for a name
    \item A dictionary should now store the grade of each name
    \begin{itemize}
        \item The name is a key, the grade the value
    \end{itemize}
    \item As soon as the user enters "exit" the program should print the grades of all names before it exits
\end{itemize}

\framebreak

\begin{itemize}
    \item Change your code in such a way that for each name it additionally outputs if the user has passed or failed
    \item Mark >= 4 -> passed
    \item Mark < 4 -> failed
\end{itemize}

\framebreak

\begin{itemize}
    \item Change your in such a way that the application outputs the average grade before it exits
\end{itemize}


\end{frame}

\begin{frame}[c, fragile]\frametitle{Exercise 10}
    
\begin{itemize}
\item  Write an application that generates a random number between 1
and 100
\item  import random
\item  random.randrange(min, max)
\item  The user makes a guess and enters a number. If the number is incorrect, the program outputs whether the entered number was too small or too large and allows the user to guess again.
\item  The application quits when the correct number is guessed
\item  The application should output how many user attempts have been
made before it quits
\end{itemize}

\end{frame}

\begin{frame}[c, fragile]\frametitle{Exercise 11}
    
    \begin{itemize}
        \item Implement the opposite of Task 10 so that the user thinks of a number and makes the computer guess
        \item The user provides feedback on whether the number is too high, too small, or correct
        \item $<$ (too low)
        \item $>$ (too high)
        \item = (correct)
        \item How many steps does the computer need?
    \end{itemize}
\end{frame}


\begin{frame}[c, fragile, allowframebreaks]\frametitle{Exercise 12}
    
\begin{itemize}
\item Write an application which repeatedly asks for a name and phone
number until the user enters ``exit''
\item Each name/telephone number pair should be stored as an entry in a dictionary
\begin{itemize}
    \item The names are the keys of the dictionary
    \item The telephone numbers are the values of the dictionary
\end{itemize}
\item As soon as the user enters ``exit'', create a JSON string of the dictionary using the json.dumps() function and store the string in a file called address\_book.txt
\end{itemize}

\framebreak

\begin{itemize}
\item Extend your application so that it reads the address\_book.txt file when it starts
\item Convert the JSON text into a dictionary again
\end{itemize}

{
\mdseries
\setbeamercolor{normal text}{fg=black}
\usebeamercolor[fg]{normal text}
\begin{lstlisting}
import json
address_book_file = open('address_book.txt', 'r')
address_book_dict = json.load(address_book_file)
\end{lstlisting}
}

\begin{itemize}
\item Ask the user if he wants to add more names or not
\item Let the user search for names in the dictionary and print out the according phone number
\end{itemize}

\end{frame}
}

\end{document}
